\subsection{Bilinear forms associated with a Hermitian form.}

We assume again here that the ring A is the quadratic extension \(\mathrm{A} = \mathrm{K}(i)\) (where \(i^{2} = -1\) ) of a commutative ring K, and that J is the automorphism \(\lambda + i\mu \to \lambda - i\mu\) of A ( \(\lambda \in K, \mu \in K\) ). Let E and \(E'\) be two A-modules, \(E_{0}\) and \(E_{0}'\) be the underlying K-modules of E and \(E'\) , j and \(j'\) the automorphisms \(x \to ix\) and \(x' \to ix'\) of E and \(E'\) (cf.~no. 2). A K-bilinear form f on \(E_{0} \times E_{0}'\) will be said to be invariant under j and \(j'\) if one has

\[
f (j (x), j ^ {\prime} (x ^ {\prime})) = f (x, x ^ {\prime})\tag{5}
\]

for \(x \in \mathrm{E}_0\) and \(x' \in \mathrm{E}_0'\) . Replacing \(x\) by \(j(x)\) , one sees that this condition is equivalent to

\[
f (x, j ^ {\prime} (x ^ {\prime})) = - f (j (x), x ^ {\prime})\tag{6}
\]

for all \(x \in \mathbf{E}_0\) and \(x' \in \mathbf{E}_0'\) .

PROPOSITION 1. --- Let \(\Phi_{1}\) (resp. \(\Phi_{2}\) ) be a K-bilinear form on \(\mathbf{E}_{0} \times \mathbf{E}_{0}^{\prime}\) , invariant under \(j\) and \(j^{\prime}\) . The map which associates to \(\Phi_{1}\) (resp. \(\Phi_{2}\) ) the map \(\Phi\) of \(\mathbf{E} \times \mathbf{E}^{\prime}\) into A defined by

\begin{enumerate}
\def\labelenumi{(\arabic{enumi})}
\setcounter{enumi}{6}
\tightlist
\item
\end{enumerate}

\[
\Phi (x, x ^ {\prime}) = \Phi_ {1} (x, x ^ {\prime}) + i \Phi_ {1} (x, j ^ {\prime} (x ^ {\prime}))\tag{8}
\]

\[
(\text { resp. } \Phi (x, x ^ {\prime}) = - \Phi_ {2} (x, j ^ {\prime} (x ^ {\prime})) + i \Phi_ {2} (x, x ^ {\prime}))
\]

\((x \in \mathrm{E}, x' \in \mathrm{E})\) , is an isomorphism of the K-vector space of K-bilinear forms on \(\mathrm{E}_0 \times \mathrm{E}_0'\) invariant under \(j\) and \(j'\) onto the K-vector space of sesquilinear forms on \(\mathrm{E} \times \mathrm{E}'\) . Suppose moreover \(\mathrm{E} = \mathrm{E}'\) ; for \(\Phi\) to be Hermitian, it is necessary and sufficient that \(\Phi_1\) be symmetric (resp. that \(\Phi_2\) be antisymmetric) (cf.~exerc. 4).

Indeed every map \(\Phi\) of E × E' into A can be written, in one and only one way, in the form \(\Phi = \Phi_1 + i\Phi_2\) , where \(\Phi_1\) and \(\Phi_2\) are maps of E × E' into K. For the partial map \(x \to \Phi(x, x')\) to be A-linear, it is necessary and sufficient, according to formula (3) (resp. (4)) of no. 2, that \(\Phi_1\) (resp. \(\Phi_2\) ) be K-linear in \(x\) and that one has

\begin{enumerate}
\def\labelenumi{(\arabic{enumi})}
\setcounter{enumi}{8}
\tightlist
\item
\end{enumerate}

\[
\Phi (x, x ^ {\prime}) = \Phi_ {1} (x, x ^ {\prime}) - i \Phi_ {1} (j (x), x ^ {\prime})\tag{10}
\]

\[
(\text { resp. } \Phi (x, x ^ {\prime}) = \Phi_ {2} (j (x), x ^ {\prime}) + i \Phi_ {2} (x, x ^ {\prime})).
\]

Similarly, for \(\overline{\Phi(x,x^{\prime})}\) to be A-linear in \(x^{\prime}\) , it is necessary and sufficient that \(\Phi_{1}\) (resp. \(\Phi_{2}\) ) be K-linear in \(x^{\prime}\) and that one has

\begin{enumerate}
\def\labelenumi{(\arabic{enumi})}
\setcounter{enumi}{10}
\tightlist
\item
\end{enumerate}

\[
\Phi (x, x ^ {\prime}) = \Phi_ {1} (x, x ^ {\prime}) + i \Phi_ {1} (x, j ^ {\prime} (x ^ {\prime}))\tag{12}
\]

\[
(\text { resp. } \Phi (x, x ^ {\prime}) = - \Phi_ {2} (x, j ^ {\prime} (x ^ {\prime})) + i \Phi_ {2} (x, x ^ {\prime})).
\]

It follows immediately that, for \(\Phi\) to be sesquilinear, it is necessary and sufficient that it be written in one or the other of the forms (9) and (11) (resp. (10) and (12)) with \(\Phi_{1}\) (resp. \(\Phi_{2}\) ) K-bilinear invariant under \(j\) and \(j'\) .

It follows from this that, for a sesquilinear form \(\Phi = \Phi_{1} + i\Phi_{2}\) to be zero, it is necessary and sufficient that \(\Phi_{1}\) (resp. \(\Phi_{2}\) ) be zero. Now, if \(E = E'\) , one has \(\Phi(y, x) = \Phi_{1}(y, x) + i\Phi_{2}(y, x)\) and \(\overline{\Phi(x, y)} = \Phi_{1}(x, y) - i\Phi_{2}(x, y)\) for these two expressions to be equal, in other words for \(\Phi\) to be Hermitian, it is therefore necessary and sufficient that \(\Phi_{1}\) be symmetric (resp. that \(\Phi_{2}\) be antisymmetric).

Remarks. --- 1) Formulas (7) and (8) show that, if \(x \in E\) , for \(\Phi(x, x') = 0\) to hold for all \(x' \in E'\) , it is necessary and sufficient that \(\Phi_1(x, x') = 0\) (resp. \(\Phi_2(x, x') = 0\) ) for all \(x' \in E'\) .

\begin{enumerate}
\def\labelenumi{\arabic{enumi})}
\setcounter{enumi}{1}
\tightlist
\item
  The adjoint of an endomorphism \(u\) of E with respect to \(\Phi\) ( \(\S 1\) , no. 8) is the same as the adjoint of \(u\) (considered as an endomorphism of \(\mathrm{E}_{0}\) ) with respect to \(\Phi_{1}\) (resp. \(\Phi_{2}\) ).
\end{enumerate}

